\id{МРНТИ 06.39.31}{}

\begin{header}
\swa{Ш.А.~Смагулова, С.З.~Кажимуратова, Г.К.~Койшинова, А.А.~Жетписбаева, С.Т.~Жакупова, М.Д.~Ахметова}{РАЗВИТИЕ ЦИФРОВОЙ ТРАНСФОРМАЦИИ ГОСУДАРСТВЕННОГО УПРАВЛЕНИЯ ЭКОНОМИКИ КАЗАХСТАНА}

\tsp{1}Ш.А.~Смагулова,
\tsp{2}С.З.~Кажимуратова\envelope,
\tsp{3}Г.К.~Койшинова,
\tsp{4}А.А.~Жетписбаева,
\tsp{5}С.Т.~Жакупова,
\tsp{6}М.Д.~Ахметова
\end{header}

\begin{affil}
\tsp{1,5}Университет Международного Бизнеса им. Кенжегали Сагадиева, Алматы, Казахстан,

\tsp{2,6}НАО Университет Нархоз, Алматы, Казахстан,

\tsp{3}Восточно-Казахстанский университет им. С.Аманжолова, Усть-Каменогорск, Казахстан,

\tsp{4}Казахский Национальный педагогический Университет им. Абая, Алматы, Казахстан

\corrauthor{Корреспондент-автор: saltanat.kazhimuratova@narxoz.kz}
\end{affil}

Цифровая трансформация государственного управления экономики направлена
на активный переход органов власти на информационно-коммуникационные
технологии в национальном хозяйстве в рамках внедрения и использования
инновационных электронных платформ для эффективного взаимодействия с
населением и представителями бизнес-среды.

Цель исследования заключается в обосновании социально-экономических
факторов и представлении рекомендаций для совершенствования
государственного управления цифровой трансформации экономического
развития РК при помощи построения многофакторного регрессионного
анализа. При написании статьи были применены методы - системного и
логического обобщения, критической оценки библиографии,
экономико-математического моделирования и прогнозирования, анализа и
статистики. Полученные результаты и выводы - осуществлен критический
обзор первоисточников с точки зрения состояния и анализа реформ
государственного управления экономическим развитием, реализован
социально-экономический анализ электронного состояния госулуг, величины
госслужащих, цифровых сервисов, проведено эконометрическое моделирование
и даны прогнозы влияния цифровой инфраструктуры на уровень роста темпов
ВВП Казахстана, приведены проблемы и представлены предложения по
улучшению организации цифровой трансформации государственного
регулирования экономического роста в РК.

{\bfseries Ключевые слова:} цифровые платформы, инвестиции государства и
предприятий в ИИ, государственные электронные услуги, цифровая
трансформация, управление экономикой.

\begin{header}
ҚАЗАҚСТАН ЭКОНОМИКАСЫН МЕМЛЕКЕТТІК БАСҚАРУДЫҢ ЦИФРЛЫҚ ТРАНСФОРМАЦИЯСЫН ДАМЫТУ

\tsp{1}Ш.А.~Смагулова,
\tsp{2}С.З.~Қажымұратова\envelope,
\tsp{3}Г.К.~Қойшинова,
\tsp{4}А.А.~Жетпісбаева,
\tsp{5}С.Т.~Жакупова,
\tsp{6}М.Д.~Ахметова
\end{header}

\begin{affil}
\tsp{1,5}Кенжеғали Сағадиев атындағы Халықаралық Бизнес Университеті, Алматы, Қазақстан,

\tsp{2,6}КЕАҚ Нархоз Университеті, Алматы, Қазақстан,

\tsp{3}С. Аманжолов атындағы Шығыс Қазақстан университет, Өскемен, Қазақстан,

\tsp{4}Абай атындағы Қазақ ұлттық педагогикалық университеті, Алматы, Қазақстан,

e-mail: saltanat.kazhimuratova@narxoz.kz
\end{affil}

Мемлекеттік экономиканы басқарудың цифрлық трансформациясы халықпен және
бизнес ортасының өкілдерімен тиімді өзара әрекеттесу үшін инновациялық
электрондық платформаларды енгізу және пайдалану арқылы мемлекеттік
органдардың ұлттық экономикада ақпараттық-коммуни\-кациялық
технологияларға белсенді көшуіне бағытталған.

Зерттеудің мақсаты - көп айнымалы регрессиялық талдауды қолдана отырып,
Қазақстан Республикасындағы экономикалық дамудың цифрлық
трансформациясын мемлекеттік басқаруды жетілдіру бойынша
әлеуметтік-экономикалық факторларды негіздеу және ұсыныстар ұсыну.
Мақаланы жазу кезінде келесі әдістер қолданылды: жүйелік және логикалық
жалпылау, библиографияны сыни бағалау, экономикалық және математикалық
модельдеу және болжау, талдау және статистика. Алынған нәтижелер мен
қорытындыларға мемлекеттік экономикалық дамуды басқару реформаларының
жай-күйі мен талдауы тұрғысынан бастапқы дереккөздерге сыни шолу;
мемлекеттік қызметтердің электрондық жағдайының, мемлекеттік
қызметшілердің санының және цифрлық қызметтердің әлеуметтік-экономикалық
талдауы; цифрлық инфрақұрылымның Қазақстандағы ЖІӨ өсу деңгейіне әсерін
эконометрикалық модельдеу және болжамдар; мәселелер талқыланып, ҚР
экономикалық өсуді мемлекеттік реттеудің цифрлық трансформациясын
ұйымдастыруды жетілдіру бойынша ұсыныстар ұсынылған.

{\bfseries Түйін сөздер:} цифрлық платформалар, үкіметтік және кәсіпорындық
инвестициялар, мемлекеттік электрондық қызметтер, цифрлық трансформация,
экономикалық басқару.

\begin{header}
THE DEVELOPMENT OF DIGITAL TRANSFORMATION OF PUBLIC ADMINISTRATION OF THE ECONOMY OF KAZAKHSTAN

\tsp{1}Sh.~Smagulova,
\tsp{2}S.~Kazhimuratova\envelope,
\tsp{3}G.~Koishinova,
\tsp{4}A.~Zhetpisbayeva,
\tsp{5}S.~Zhakupova,
\tsp{6}M.~Akhmetova
\end{header}

\begin{affil}
\tsp{1,5}Kenzhegali Sagadiyev University of International Business, Almaty, Kazakhstan,

\tsp{2,6}Narxoz University, Almaty, Kazakhstan,

\tsp{3}East Kazakhstan University named after S. Amanzholov, Ust-Kamenogorsk, Kazakhstan,

\tsp{4}Abai Kazakh Pedagogical University, Almaty, Kazakhstan,

e-mail: saltanat.kazhimuratova@narxoz.kz
\end{affil}

The digital transformation of public economic management is aimed at the
active transition of govern\-ment bodies to information and communication
technologies in the national economy through the intro\-duction and use of
innovative electronic platforms for effective interaction with the
population and repre\-sentatives of the business environment.

The purpose of the study is to substantiate socio-economic factors and
present recommendations for improving the public administration of the
digital transformation of economic development in the Republic of
Kazakhstan through a multivariate regression analysis. When writing the
article, the following methods were used: systemic and logical
generalization, critical assessment of bibliography, economic and
mathe\-matical modeling and forecasting, analysis and statistics. The
results and conclusions obtained include a critical review of primary
sources from the point of view of the state and analysis of public
economic development management reforms; a socio-economic analysis of
the electronic state of public services, the number of civil servants,
and digital services; econometric modeling and forecasts of the impact
of digital infrastructure on the level of GDP growth in Kazakhstan;
problems are discussed and proposals for improving the organization of
the digital transformation of state regulation of economic growth in the
Republic of Kazakhstan are presented.

{\bfseries Keywords:} digital platforms, government and enterprise
investments in AI, government e-services, digital transformation,
economic management.

\begin{multicols}{2}
{\bfseries Введение.} В феврале 2026г. А.Гутерриш - Генеральный секретарь
ООН отметил значительную роль научных исследований и инноваций в
международном управлении ИИ\emph{.} В частности, он подчеркнул, что
инновации в области искусственного интеллекта развиваются со скоростью
света, опережая нашу коллективную способность в полной мере понимать их,
не говоря уже об управлении ими» {[}1{]}.

По данным ООН ускорение цифрового регулирования экономических процессов
привело к росту мировой торговли до уровня 7\%, что в абсолютной
величине составляет около 35 трлн. долл. США.\emph{~} При этом,
цифровизация сферы услуг приблизилась к 56\% от всего мирового экспорта
{[}2{]}.

Другие исследователи обращают внимание на аспекты устойчивости цифровых
инициатив на примере оценки муниципалитетов Исландии (Sigurjonsson et
al.2024). В частности, проведен анализ предоставления государственных
услуг. На основе комплексного подхода авторы определили, что для
оптимизации государственных услуг требуется активизировать процесс
цифровой трансформации и автоматизации государственного управления
{[}3{]}.

Согласно оценкам
«\href{https://datareportal.com/reports/digital-2026-global-overview-report}{Digital
2026 Global Overview Report}» сегодня 5,78 млрд. людей пользуются
мобильными устройствами (чуть больше 70\% от всего населения Земли). К
началу 2026 года количество пользователей интернета в мире составило
6,04 млрд. человек. Если вернуться к истории появления первого вэб-сайта
в 1991 г., то в тот период количество пользователей было приблизительно
5 млн. человек в мире {[}4{]}.

Экономика Казахстана, очевидно, имеет высокую зависимость от
углеводородного сырья. В этой связи необходимо активизировать процесс
диверсификации секторов национального хозяйства. При этом акцент следует
уделить структуре ВВП, экономическому развитию обрабатывающего сектора,
электронным услугам и цифровым технологиям.

Выделим, 2026 г. объявлен в нашей стране - Годом развития цифровизации и
ИИ с целью внедрения инноваций, обеспечения конкурентоспособности и
устойчивого экономического роста. Необходимость цифровизации управления
экономики продиктована ростом эффективности оказания госуслуг,
прозрачности государственного регулирования экономических процессов. К
тому же она стимулирует внедрение в деятельность отраслей IT, AI,
помогает улучшать и развивать финансовые технологии на базе ИИ с
использованием инновационных технологий, создавать новые рабочие места в
сфере ИКТ, укреплять макроэкономическую стабильность и устойчивость
государства.

Используя программу «VOSviewer» Smagulova et al. (2025), реализовали
библиометрическую оценку лучших международных статей по вопросам
цифровизации государственного регулирования с 1985 по 2024 годы {[}5{]}.
Авторы сделали акцент на применении цифровых элементов при осуществлении
зарубежных концепций, подходов и программ в сфере государственного
управления экономическими процессами на примере лучших мировых практик.

С целью развития цифровизации в нашей стране приняты такие важные
законодательные документы:

- Цифровой Кодекс РК, 2026 {[}6{]};

- Закон Республики Казахстан «Об искусственном интеллекте», 2025
{[}7{]};

- «Концепция развития искусственного интеллекта на 2024-2029 гг.», 2024
{[}8{]} и др.

Например, согласно Закону «Об искусственном интеллекте» - ИИ
представляется объектом информатизации и важным инструментом в
деятельности гражданина. Вместе с тем все пользователи несут полную
ответственность по соблюдению безопасности, конфиденциальности и
прозрачности в сфере работы ИИ.

В 2025 году в Казахстане организовано новое Министерство искусственного
интеллекта и цифрового развития. Предполагается, что данная
государственная структура будет позитивно воздействовать на
последовательном и системном внедрении инструментов ИИ, развитии
технологий и инновационной инфраструктуры, а также способствовать
приобретению цифровых компетенций и квалификации кадров.

Отдельное внимание требуется уделить вопросам цифровизации
государственного управления экономикой. Это предполагает повышение
эффективности и рост информационно-технологического тренда общей
экономической системы. Цифровые технологии непосредственно позволяют
ускорить принятие управленческих решений, снизить издержки
государственного управления, минимизировать погрешности менеджмента
государственных служащих.

К тому же совершенствование электронных сервисов государства дают
возможность - улучшить качество государственных услуг, уменьшить
административные барьеры, казахстанским гражданам мобильнее получать
услуги в онлайн формате, сократить очереди и сроки получения
разрешительных документов, упростить процесс регистрации бизнеса и
ведения налоговой отчётности предприятий и пр.

В настоящее время благоприятными предпосылками формирования
технологического развития и искусственного интеллекта (ИИ) в Казахстане
представляются: цифровой потенциал, возможность адаптации
законодательной базы, определенное наличие массива статданных,
налаживание контактов с ведущими технологическими компаниями, увеличение
государственной поддержки для стимулирования использования ИИ и др.

Существенной составляющей здесь также является цифровизация
государственных услуг. Определенный уровень развития получил
финтех-сектор и экосистема финансирования новых «IT-стартапов». Так,
национальный инновационный кластер «Astana Hub» включает около 2000
предприятий, которые в том числе в 2025 г. обеспечили экспорт «IT-услуг»
на сумму в размере 1 млрд долл. США.

Значительным позитивным условием организации новых ИКТ представляется
запуск 2-х крупных суперкомпьютеров «Alem.Cloud» и «Al-Farabium» с
высокой производительностью и применением чипов «NVIDIA».

Цифровизация госуправления стимулирует развитие: обрабатывающих
отраслей, внедрение промышленных инновационных технологий, «e-commerce»,
«fintech», автоматизации разных социальных выплат и др.

Электронные и цифровые платформы способствуют понятной деятельности
государства, развитию конкурентной среды, объективно принимать
обоснованные решения и позволяют обществу мониторить управленческие
процессы.

В конечном результате трансформация цифровизации экономической системы
государственного управления позволит работать оперативнее, доступнее и
прозрачнее, а также шире укреплять доверие общества к госорганам и
улучшать эффективность реализации социальной политики.

Однако, серьезными недостатками развития цифровых технологий в
республике сегодня являются - слабая степень развитости экосистемы
инфраструктуры, низкие возможности потенциала инноваций, недостаточное
количество больших технологических компаний, низкие уровни
государственных и частных расходов на опытно-конструкторские и
научно-исследовательские мероприятия, слабые возможности потенциала
человеческого капитала и продвижения информационной культуры граждан,
недостаточность внедрения интеллектуальных этических норм.

{\bfseries Методы и материалы.} В научном исследовании Tutar
(2025) определяется взаимодействие прозрачных отношений
государства, политической активности и общественности в условиях
цифровизации {[}9{]}. С помощью эконометрического моделирования и
панельной регрессии были получены такие выводы. В частности,
деятельность в социальных сетях стимулирует общественность для активного
участия в работе электронного правительства. В свою очередь, это
способствует продвижению и укреплению цифровой демократии в системе
госуправления.

В этой работе будет построена множественная регрессионная модель с
несколькими переменными:

\begin{equation}
Y = f(X_{1}, X_{2}, X_{3}, \ldots, X_{m}) + \epsilon
\end{equation}

Представим уравнение модели множественной регрессии, которая будет
использована в нашем исследовании:

\begin{equation}
Y = a_{0} + a_{1}X_{1} + a_{2}X_{2} + a_{3}X_{3} + \cdots + a_{m}X_{m} + \epsilon
\end{equation}

где:

\(Y\) - зависимый показатель,

\(X_{1}, X_{2}, \ldots, X_{m}\)- независимые факторы,

\(a_0\) - свободный член модели,

\(a_{1}, a_{2}, \ldots, a_{m}\) - коэффициенты регрессии, характеризующие влияние
определенного независимого показателя,

\(\varepsilon\) - случайная ошибка.

Здесь определенный коэффициент \(a_i\) демонстрирует, как изменится
зависимый показатель \(Y\) при повышении некоторой переменной
\(X_i\) на одну единицу, если все другие переменные будут при этом
неизменными. По мнению
\href{https://journals.sagepub.com/doi/10.1177/00208523251333541\#con}{Borriello}
(2025), цифровизация государственного управления на примере Италии
демонстрирует важное значение для решения конфликтов в принятии
определенных решений {[}10{]}. Ученый в своем исследовании с помощью
моделирования экономических процессов на базе прикладной программы
«NVivo» доказывает, что электронные и цифровые режимы положительно
воздействуют на реформировании государственного сектора, обеспечения
национальной безопасности и осуществления объективного мониторинга.

Цель представленного эконометрического моделирования - построить модель
с определенным количеством нескольких факторов одновременно, чтобы
улучшить качество анализа и глубину краткосрочного прогноза.~При этом,
определяется воздействие каждого фактора и общее их влияние на зависимую
переменную. Алгоритм многофакторного анализа включает такие этапы:
проводится сбор, очистка и отбор независимых факторов, влияющих на
исследуемую переменную, анализируется первоначальная информация,
осуществляется связь между независимыми факторами и результативным
индикатором, строится полученное уравнение, реализуется расчет
переменных взаимодействия множественного корреляционного анализа,
выводится статанализ итоговых результатов, строится прогноз.

Одной из значительных проблем в системе государственного управления
представляется внедрение и использование инновационных ИКТ-решений
(Koerich et al.2025). Авторы провели интервью и собрали данные по
оценке факторов, влияющих на использование цифровых технологий {[}11{]}.
Так, используя методику контент-анализа были выявлены следующие факторы:
социально-политические, развитие информационной инфраструктуры, рост
уровня квалификации государственных служащих, качество человеческого
капитала, цифровизация госуслуг и др.

Индикаторы линейного уравнения многофакторной регрессии, как правило,
проводятся методом наименьших квадратов (МНК):

\begin{equation}
\sum_{i} (y_{i} - \hat{y}_{i})^{2} \to \min
\end{equation}

Для эконометрического моделирования в этой статье были поставлены
гипотезы:

- на уровень цифровой трансформации управления экономическим развитием
РК наиболее сильное воздействие на современном этапе оказывают -
Инвестиции государства и предприятий в информационно-коммуникационные
технологии;

- по прогнозам в краткосрочной перспективе номинальный ВВП ежегодно
будет расти в коридоре 10--18 процентов за счет внедрения цифровых
факторов в Казахстане.

В качестве инструмента используется метод наименьших квадратов (OLS)
анализа. Спецификация модели выглядит так:
\end{multicols}
\vspace{-0.5cm}
\begin{equation}
Y_{t} = \beta_{0} + \beta_{1}X1_{t} + \beta_{2}X2_{t} + \beta_{3}X3_{t} + \beta_{4}X4_{t} + \beta_{5}X5_{t} + \beta_{6}X6_{t} + \epsilon_{t}
\end{equation}

\begin{multicols}{2}
Размер выборки составил 16 лет (2010--2025 гг.). Доверительный интервал
прогноза в 5\% подтверждает статистическую достоверность полученных
расчетов.

\begin{equation}
\hat{Y}_{t+h} \pm t_{\alpha/2} \cdot SE(\hat{Y})
\end{equation}

где:

\(t_{\alpha/2}\) - показатель критического распределения критерия
Стьюдента,

\(SE(\hat{Y})\) - стандартная ошибка прогноза.

Цифровизация государственного управления на основе подхода открытых
инноваций играет непосредственную роль в совершенствовании развития
корпоративной цифровизации, научных знаний и накоплении информационного
капитала. Так, по мнению китайских ученых (Yang et al.2025) в рамках
математического моделирования было обоснованно, что существует прямая
связь между цифровыми платформами Правительства, используемыми
технологиями и инновационной деятельности предприятий {[}12{]}.

Для построения и получения результатов эконометрической модели был
использован прикладной пакет «Gretl». В данном пакете содержатся готовые
встроенные этапы с целью оценки временных рядов. К примеру, в этой
программе существует алгоритм анализа индикаторов на основе -- «МНК»
(т.е. метода наименьших квадратов), метода моментов и т.д. Анализ
временных рядов включает -- «ARMA» (статистическая модель -
авторегрессия скользящего среднего), «ARIMA» (попытка реализовать
результат будущего показателя на основе его прошлого состояния -
авторегрессия комплексного скользящего среднего), «VAR» векторная
авторегрессия и пр.

В качестве информационной базы экономико-математического моделирования
были использованы материалы, справки, отчеты и информационные бюллетени
- Всемирного Банка, Концепций, национальных проектов и госпрограмм
Казахстана, Национального бюро по статистике и пр.

{\bfseries Обсуждение и результаты.} По результатам деятельности
цифровизации экономики РК, международный аналитический центр «Oxford
Insights» сообщает, что наша республика заняла почетное 60-е место среди
исследуемых 195 стран в рейтинге «Government AI Readiness Index 2025»
{[}13{]}. По сравнению с прошлым периодом, рейтинг возрос на 16 мест
(ранее 76-я строчка), что демонстрирует достаточное лидерство Казахстана
среди государств Центральной Азии в вопросах подготовленности к
применению и управлению ИИ в национальной экономике.

В настоящее время использование цифровых технологий в электронном
правительстве вызывает огромное внимание к intelligent decision support
systems (IDSS) {[}14{]}. С целью активизации внедрения цифровых решений
в госуправлении ученые Chen et al. (2025). рекомендуют использовать
«Integrated AI-based Decision Support System» (IA-DSS-EG) на основе ИИ.
Отмеченная система состоит из таких 5 модулей. К ним можно отнести:
общая классификация административных функций, прогнозирование налогов на
развитие экономики, степень доверия общества, соблюдения правовой
совместимости и повышение квалификации госслужащих.

Отметим, рейтинг «Government AI Readiness Index» включает оценку
возможностей правительства рационально применять ИИ в экономике,
государственном управлении и др. Значительные достижения в рамках
отмеченного рейтинга республика показала по разделу «Public Sector
Adoption» - 73,59 балла. Такой результат говорит о разносторонних
возможностях электронного правительства, цифровизации госуслуг,
организации информационных сервисов ИИ для юридических и физических лиц,
роста финансирования бизнеса для внедрения цифровых решений.

На основе исследований Казахстан по качественной степени цифрового
управления уверенно вошел в~10-ку основных лидеров глобального рейтинга
«OS»I (Online Services Index) в 2025 году {[}15{]}. В частности, по
итогам 2026 года свыше 94\% госуслуг предлагаются в электронном формате
(рисунок 1).

Из данных рисунка 1 наблюдается постепенный рост зарплаты. При этом
существует обратная зависимость между некоторым снижением количества
госслужащих и уровнем зарплаты. Так, величина зарплаты возросла со
125000 до 420000 тенге или в 3,4 раза за последние 10 лет. А количество
чиновников со 100000 упало до 84000 человек (или на 16000 чел.) с 2016г.
до 2025г. Другими словами, оптимизация численности государственных
сотрудников привела к учеличению эффективности госаппарата и оплаты
труда.
\end{multicols}

{\bfseries Рис.1 - Состояние развития электронных госуслуг и госслужащих в
РК, 2016--2025}

\emph{Примечание: Составлено на основе источников {[}16 - 18{]}}

\begin{multicols}{2}
Dei, 2025 сформировал важные профессиональные компетенции госслужащих,
которые требуются для работы с электронными технологиями в сфере ИИ
{[}19{]}. Особо выделяется, что цифровизация системы местного
самоуправления дает возможность применять методику ИИ для принятия
перспективных государственных решений.

Считаем, что одним из важных факторов эффективности состава госаппарата
и роста зарплаты по результатам оценки деятельности чиновников
представляется цифровизация и использование ИИ при принятии
государственных решений.

Коррупция представляется серьезной проблемой для обеспечения
прозрачности в решении вопросов государственного регулирования
экономических процессов.
\href{https://www.researchgate.net/scientific-contributions/Darusalam-Darusalam-2146395918?_tp=eyJjb250ZXh0Ijp7ImZpcnN0UGFnZSI6ImNpdGF0aW9uRG93bmxvYWQiLCJwYWdlIjoicHVibGljYXRpb24ifX0}{Darusalam}
et al. (2024) провели качественное исследование в виде интервью с
государственными служащими и представителями населения, чтобы вскрыть
угрозы автоматизации и цифровизации административных процессов в органах
государственной службы {[}20{]}. Специалисты установили, что следует
внедрять электронные технологии с целью снижения сговора и возможности
злоупотребление среди государственных лиц.

В 2025 году пользователям предложены более 1200 цифровых сервисов
посредством платформ -- «eGov.kz» и «eGov Mobile». При этом объем
предоставленных госуслуг в РК составил свыше 520 млн. единиц, рост
составляет около 15 млн ед. по сравнению с 2024 годом.

Здесь огромную поддержку указанным платформам оказывает виртуальный
помощник~- «eGov AI» особенно в части ответов на~присланные вопросы о
видах, сроках и объемах~онлайн-госуслуг. Этот помощник значительно
сократил время на ожидание вопросов граждан и перевел ручную работу на
-автоматическую в~госорганах, а также снизил нагрузку и расходы
на~работу операторов. Портал «eGov» включает регистрацию услуг
предпринимателям, медуслугам, подачи заявлений и др. необходимых
документов. Электронный сервис «eGov Mobile» использует элементы
«push-уведомлений», процессов биометрической идентификации и~пр.

В работе Urhibo et al. (2024) проведена оценка государственных услуг в
американской сфере госуправления. Реализованный опрос 146 респондентов
выявил, что существует возможность внедрения электронных и цифровых
технологий в обеспечении государственных услуг {[}21{]}. В этой связи
необходимо усиленно вкладывать инвестиции в развитии электронной
инфраструктуры, расширения цифровой грамотности среди населения и
повышения профессионализма государственных кадров.

В результате идет активный процесс комплексной интеграции ИИ в разные
отрасли экономики в: государственное управление, энергетику, АПК,
здравоохранение и др.

Значит цифровая модернизация включает не только сам доступ к электронным
услугам и сервисам, но и критерии управления статданными.

Активизация технологической модернизации напрямую связана с созданием
стартапов, которые привлекли более 300~млн долл. США венчурных
и государственных инвестиций. В частности, крупный местный
акелератор «Astana Hub» объединил свыше 1350 стартапов из~таких стран:
США, Англия, Голландия, ФРГ, Россия и др. (более 30 государств), общий
доход равен 1,3~трлн. тенге.

Стоит выделить, что в республике был открыт в 2025 году крупный
венчурный фонд «Qazaqstan Venture Group», величина которого составляет 1
млрд. долл. США. Основная деятельность данного фонда ориентирована на
финансирование ИИ-стартапов и развитие IT-инфраструктуры.

Цифровая модернизация воздействует также и на рынок труда. Уже
зарегистрировано около 18,5 тыс. IT-компаний, где осуществляют трудовую
деятельность более 195 тыс. человек. Разработанная и реализуемая
программа «Tech Orda» предполагает подготовить к 2029 году 20 тыс. новых
IT-специалистов и инженеров ИИ при поддержке государства.

Следовательно, все вышесказанное демонстрирует стратегию мобильной
направленности, финансирования, модернизации, управления цифровой
системы и стимулирования роста экономики страны в Казахстане.

Для научно-практического исследования влияния цифровых факторов на
экономическое развитие Республики Казахстан была реализована
множественная корреляция (таблица 1) между номинальным ВВП (\emph{Y})
страны и важными переменными цифровой инфраструктуры (\emph{Х1}),
инвестиций (\emph{Х2}), производительностью человеческого капитала
(\emph{Х4, Х5}) и гражданами, владеющими информационными ресурсами
(\emph{Х3, Х6}).
\end{multicols}

\tcap{Таблица 1 -- Исходные статданные для построения модели}
\begin{longtblr}[
  label = none,
  entry = none,
]{
  width = \linewidth,
  colspec = {Q[42]Q[100]Q[100]Q[100]Q[100]Q[100]Q[100]Q[100]},
  cells = {c},
  cells = {font = \footnotesize},
  row{2} = {font=\bfseries\footnotesize},
  row{19} = {font=\itshape},
  cell{1}{1} = {r=2}{font=\bfseries\footnotesize},
  cell{1}{5} = {font=\bfseries\footnotesize},
  cell{1}{6} = {font=\bfseries\footnotesize},
  cell{1}{8} = {font=\bfseries\footnotesize},
  cell{19}{1} = {c=8}{0.935\linewidth},
  hlines,
  vlines,
}
Годы & {\textbf{Номиналь}\\\textbf{-ный}\\\textbf{ВВП РК,}\\\textbf{Млрд.долл. США}} & {\textbf{Инвестиции государства}\\\textbf{в цифровые технологии}\\\textbf{в РК}\\\textbf{(млрд .тг)}} & {\textbf{Инвестиции предприятий}\\\textbf{в цифровые технологии в РК (млрд .тг)}} & Уровень цифровой грамотности населения РК (\%) & Индекс человеческого капитала РК (0–1) & {\textbf{Производи-тельность труда РК (Информация}\\\textbf{и связь) (тыс. тг.)}} & Доля населения, пользующе-гося интернетом (\%)\\
 & Y & Х1 & Х2 & Х3 & Х4 & Х5 & Х6\\
2 010 & 148,00 & 30,00 & 120,00 & 62,00 & 0,77 & 54 450,00 & 31,62\\
2 011 & 192,00 & 35,00 & 150,00 & 65,00 & 0,77 & 45 788,40 & 50,60\\
2 012 & 208,00 & 40,00 & 180,00 & 67,00 & 0,78 & 51 083,00 & 61,90\\
2 013 & 236,00 & 45,00 & 210,00 & 69,00 & 0,79 & 63 687,80 & 63,30\\
2 014 & 221,00 & 55,00 & 240,00 & 72,00 & 0,79 & 61 157,10 & 63,90\\
2 015 & 184,00 & 60,00 & 260,00 & 73,00 & 0,80 & 66 712,20 & 72,90\\
2 016 & 137,00 & 70,00 & 290,00 & 74,00 & 0,80 & 60 581,00 & 95,00\\
2 017 & 166,00 & 85,00 & 350,00 & 76,00 & 0,802 & 68 960,60 & 94,00\\
2 018 & 179,00 & 100,00 & 420,00 & 79,00 & 0,804 & 71 981,10 & 93,00\\
2 019 & 181,00 & 120,00 & 480,00 & 82,00 & 0,810 & 80 287,30 & 91,00\\
2 020 & 171,00 & 140,00 & 550,00 & 84,00 & 0,83 & 90 530,80 & 89,00\\
2 021 & 197,00 & 160,00 & 650,00 & 87,00 & 0,82 & 103 061,70 & 87,00\\
2 022 & 225,00 & 180,00 & 820,00 & 88,00 & 0,83 & 130 940,80 & 84,00\\
2 023 & 262,00 & 210,00 & 980,00 & 90,00 & 0,84 & 136 907,00 & 82,00\\
2 024 & 291,00 & 250,00 & 1 150,00 & 91,00 & 0,64 & 151 898,70 & 79,00\\
2 025 & 300,00 & 320,00 & 1 250,00 & 94,00 & 0,65 & 103 179,00 & 76,00\\
\textit{Примечание: Составлено на основе источников 22 - 26} &  &  &  &  &  &  & 
\end{longtblr}

В нижеприведенной таблице 2 представлены коэффициенты корреляции между
зависимой переменной \emph{Y} (номинальный ВВП РК) и различными
независимыми индикаторами, а также между самими независимыми факторами.

\tcap{Таблица 2 - Коэффициенты корреляции}
\begin{longtblr}[
  label = none,
  entry = none,
]{
  cells = {c},
  cells = {font = \small},
  row{1} = {font=\bfseries\itshape\small},
  row{8} = {font=\itshape\small},
  cell{1}{1} = {font=\bfseries\small},
  cell{2}{1} = {font=\itshape\small},
  cell{3}{1} = {font=\itshape\small},
  cell{4}{1} = {font=\itshape\small},
  cell{5}{1} = {font=\itshape\small},
  cell{6}{1} = {font=\itshape\small},
  cell{7}{1} = {font=\itshape\small},
  cell{8}{1} = {c=7}{},
  hlines,
  vlines,
}
Независимые факторы & Y & Х1 & Х2 & Х3 & Х4 & Х5\\
\textit{Х1} & 0,734 &  &  &  &  & \\
\textit{Х2} & 0,762 & 0,991 &  &  &  & \\
\textit{Х3} & 0,595 & 0,939 & 0,939 &  &  & \\
\textit{Х4} & -0,611 & -0,501 & -0,482 & -0,240 &  & \\
\textit{Х5} & 0,658 & 0,851 & 0,905 & 0,884 & -0,247 & \\
\textit{Х6} & -0,021 & 0,412 & 0,404 & 0,648 & 0,220 & 0,408\\
\textit{Примечание: построено на основе программы \textbf{«Gretl»}} &  &  &  &  &  & 
\end{longtblr}

Приведенный множественный анализ показывает, что основными движущими
силами увеличения номинального ВВП представляются инвестиции в цифровые
технологии и производительность труда. Государственные инвестиции
(\emph{Х1}) (r = 0,734) и частные инвестиции (\emph{Х2}) (r = 0,762)
показывают сильную положительную cвязь с уровнем ВВП. Соответственно,
это значительно подтверждает важность как бюджетного, так и
предпринимательского финансирования цифровизации экономики РК.

Производительность труда (\emph{Х5)} имеет положительную корреляцию (r =
0,658) и демонстрирует связь между показателями инвестиций и
экономическим развитием. В то же время степень цифровой грамотности
населения (\emph{Х3}) показывает достаточно умеренную положительную
связь с ВВП (r = 0,595). Это подчеркивает определенное влияние
инновационной деятельности граждан на зависимую переменную \emph{(Y)}.

Однако, индекс человеческого капитала (\emph{Х4}) имеет отрицательную
корреляцию (r = -0,611), что, по нашему мнению, объясняется некоторой
временной задержкой его воздействия (т.е. проявится вероятно через
длительное время). Доля населения, пользующегося интернетом (\emph{Х6}),
в настоящее время имеет очень слабую отрицательную корреляцию (r =
-0,021). Это говорит о том, что индикатор пока не оказывает
значительного влияния на ВВП.

Определим взаимосвязь между зависимой переменной (\emph{Y}) и
независимыми переменными (\emph{Х1-Х6}). В частности, существует почти
функциональная взаимосвязь между государственными и частными
инвестициями (r = 0,991), что указывает на синхронизацию цифрового
развития. Другими словами, здесь синхронизация подразумевает, что данные
два фактора одновременно участвуют в продвижении и управлении цифрового
развития экономики РК. ~

Инвестиции тесно связаны с цифровой грамотностью (r = 0,939), а цифровая
грамотность тесно связана с производительностью труда (r = 0,884). Объем
производства тесно связан с инвестициями (r = 0,905), что подтверждает
характер системы цифровой трансформации. Выделим, что цифровая
трансформация~в совокупности предполагает процедуру перехода
государственных органов власти и внедрение новых
информационно-коммуникационных технологий в различные сектора экономики
с позиции роста эффективности управления. Вместе с тем происходит
перестройка операционной модели государства, введение новых электронных
платформ для обоюдовыгодных контактов Правительства с населением и
бизнесом, а также создание прозрачности управленческих решений
государства на всех уровнях.~Значит в стране в рамках цифровой
трансформации осуществляются процессы «оцифровки» государственных услуг,
при этом государственные органы ориентированы на
клиентоориентированность, конкурентоспособность и устойчивость отраслей
национальной экономики к изменениям рынка.~

Индекс человеческого капитала имеет слабые связи с другими показателями
(- 0,247 до -0,220), что указывает на то, что краткосрочное влияние
невелико. Этот Индекс (Human Capital Index), в том числе
представляет сведения об объёме и качестве образования.~В связи с тем,
что владение цифровыми навыками трудовых ресурсов еще находится на
начальном уровне, то поэтому не видно тесной связи с цифровой
трансформацией на уровне экономического развития Казахстана.

Приведенная статмодель была построена методом наименьших квадратов для
возможного решения вопросов с наличием автокорреляции и
гетероскедастичности. Для оценки значимости коэффициентов использовались
показатели «T-Value» и «P-Value» (показывает статистическую
значимость переменных, где значимость считается тогда, когда «p-значение
меньше 0.05»), а для выявления мультиколлинеарности, чтобы
исключать или совместно объединять коррелирующие показатели -- Variance
Inflation Factor (VIF).

Уравнение регрессии нашей модели выводится по данным таблицы
3. Следовательно, наше уравнение регрессии имеет следующий вид:

\begin{equation}
Y = -229 - 4,62X_{1} + 1,103X_{2} + 14,15X_{3} - 274X_{4} - 0,00329X_{5} - 2,547X_{6}
\end{equation}

\tcap{Таблица 3 -- Коэффициенты}
\begin{longtblr}[
  label = none,
  entry = none,
]{
  cells = {c},
  cells = {font = \small},
  row{1} = {font=\bfseries\small},
  hlines,
  vlines,
}
Term & Coef & SE Coef & T-Value & P-Value & VIF\\
Constant & -229 & 263 & -0,87 & 0,408 & \\
Х1 & -4,62 & 1,53 & -3,01 & 0,015 & 498,57\\
Х2 & 1,103 & 0,348 & 3,17 & 0,011 & 467,66\\
Х3 & 14,15 & 6,46 & 2,19 & 0,056 & 121,73\\
Х4 & -274 & 201 & -1,36 & 0,207 & 3,83\\
Х5 & -0,00329 & 0,00123 & -2,67 & 0,026 & 46,96\\
Х6 & -2,547 & 0,942 & -2,70 & 0,024 & 8,12
\end{longtblr}

Государственные инвестиции в цифровые технологии (\emph{Х1}).
Отрицательный коэффициент корреляции и хорошая значимость (\emph{р} =
0,015) могут отражать краткосрочные и небольшие по абсолютной величине
экономические расходы на цифровизацию, такие как: перераспределение
бюджетных средств и определенные потери при введении цифровых факторов
для создания инфраструктуры.

Инвестиции бизнеса в цифровые технологии (\emph{Х2}). Положительное и
значительное влияние (\emph{р} = 0,011) подтвердило хорошую
привлекательность частных инвестиций в продвижении и стимулировании
экономического роста. Это может характеризовать с более эффективным
приобретением инновационной техники и гибким использованием новых
технологий.

Цифровая грамотность населения (\emph{Х3}). Положительный эффект на
уровне 10\% значимости указывает на то, что повышение квалификации
населения в перспективе будет увеличивать экономический потенциал
Казахстана и способствовать увеличению ВВП.

Индекс человеческого капитала (\emph{Х4}) Отсутствие статистической
значимости (\emph{р} = 0,207) говорит о том, что в приведенной модели
пока определенного позитивного влияния человеческого капитала на темпы и
уровень объема ВВП не было установлено, что может быть связано с
ограниченными временными рядами.

Производительность труда (\emph{Х5}). Отрицательный коэффициент
(\emph{р} = 0,026) могут отражать недостаточные затраты на модернизацию
процесса образования и внедрение цифровых технологий.

Доля населения, пользующегося Интернетом (\emph{Х6}). Значительное
негативное влияние (\emph{р} = 0,024) указывает на адаптивный эффект,
когда доступ к интернет-ресурсам не конвертируется в экономическую
отдачу.

В целом результаты показывают, что на современном этапе развития
цифровой экономики частные инвестиции в большей степени способствуют
экономическому росту, в то время как государственные инвестиции и
эффективность инфраструктуры ИКТ пока еще сопровождаются краткосрочными
расходами на цифровую адаптацию.

\tcap{Таблица 4 - Краткое описание полученной модели}
\begin{longtblr}[
  label = none,
  entry = none,
]{
  cells = {c},
  cells = {font = \small},
  row{1} = {font=\bfseries\itshape\small},
  hlines,
  vlines,
}
S & R-sq & R-sq(adj) & R-sq(pred)\\
22,7345 & 86,22\% & 77,03\% & 0,00\%
\end{longtblr}

Модель демонстрирует достаточно высокую согласованность, объясняя
86,22\% изменения увеличения показателя ВВП в зависимости от выбранных
цифровых и других социально-экономических факторов (таблица 4).

Скорректированный \emph{R²} (77,03\%) показывает, что с учетом
количества переменных модель обладает достаточно высокой объясняющей
способностью с 6 независимыми факторами.

Низкий \emph{R²} (pred) = 0\% демонстрирует то, что модель не оценивает
значения ВВП за пределами исходного набора данных (ограниченные
временные ряды).

Качество модели: высокий контроль (\emph{R²}=86,22\%) позволяет сделать
вывод о важности выбранных экономических факторов для роста уровня ВВП,
в то время как низкий \emph{R²} (pred) указывает на ограниченную
предсказательную способность модели за пределами исследуемого периода.

Таким образом, можно сделать вывод, что для повышения эффективности
цифровизации в Казахстане необходимо объединить государственные и
частные инвестиции, усилить акцент на развитии цифровой грамотности и
минимизировать краткосрочные экономические издержки цифровой
трансформации.

В таблице 5 приведены результаты дисперсионного анализа (ANOVA)
множественной эконометрической модели, направленного на оценку
значимости её отдельных компонентов и формирование целостной
конфигурации модели.

\tcap{Таблица 5 - Дисперсионный анализ}
\begin{longtblr}[
  label = none,
  entry = none,
]{
  cells = {c},
  cells = {font = \small},
  row{1} = {font=\bfseries\itshape\small},
  cell{3}{1} = {font=\itshape\small},
  cell{4}{1} = {font=\itshape\small},
  cell{5}{1} = {font=\itshape\small},
  cell{6}{1} = {font=\itshape\small},
  cell{7}{1} = {font=\itshape\small},
  cell{8}{1} = {font=\itshape\small},
  hlines,
  vlines,
}
Source & DF & Adj SS & Adj MS & F-Value & P-Value\\
Regression & 6 & 29100,0 & 4850,0 & 9,38 & 0,002\\
\textit{Х1} & 1 & 4694,4 & 4694,4 & 9,08 & 0,015\\
\textit{Х2} & 1 & 5191,6 & 5191,6 & 10,04 & 0,011\\
\textit{Х3} & 1 & 2481,6 & 2481,6 & 4,80 & 0,056\\
\textit{Х4} & 1 & 957,0 & 957,0 & 1,85 & 0,207\\
\textit{Х5} & 1 & 3686,1 & 3686,1 & 7,13 & 0,026\\
\textit{Х6} & 1 & 3775,7 & 3775,7 & 7,31 & 0,024\\
Error & 9 & 4651,7 & 516,9 &  & \\
Total & 15 & 33751,8 &  &  & 
\end{longtblr}

\emph{Общая значимость модели.} Значение \emph{F-Value} = 9,38 и
\emph{P-Value} = 0,002. Результаты показывают, что регрессионная модель
в целом является статистически значимой. Это означает, что выбранные
независимые переменные имеют большое значение для объяснения изменений в
номинальном ВВП Казахстана. Таким образом, наша модель обладает высокой
объяснительной способностью.

\emph{Значимость отдельных факторов}. Инвестиции государства в цифровые
технологии в Казахстане (млрд. тг.) (\emph{Х1}). Значение
\emph{F-Value}= 9,08 и \emph{P-Value} = 0,015. Это означает, что
государственные инвестиции в цифровые технологии оказывают негативное
влияние на номинальный ВВП Казахстана. Это может быть связано с
краткосрочными затратами, необходимыми для внедрения новых технологий,
инфраструктурных проектов и цифровых инициатив. Однако, несмотря на
отрицательное значение коэффициента (-4,62), высокая значимость этого
фактора подтверждается статистическими тестами.

Инвестиции предприятий в цифровые технологии в Казахстане (млрд. тг.)
(\emph{Х2}). Таким образом, значение \emph{F-Value} = 10,04 и
\emph{P-Value} =0,011. Частные инвестиции в цифровые технологии
положительно влияют на номинальный ВВП Казахстана. Положительный
коэффициент (1,103) указывает на то, что увеличение частных инвестиций
на 1 млрд. тг. может привести к росту ВВП на 1,103 млрд. USD. Это
подтверждает важность частных инвестиций в цифровую трансформацию
экономики.

Уровень цифровой граммотности населения Казахстан (\%) (\emph{Х3}) -
значение \emph{F-Value} = 4,80 и \emph{P-Value} = 0,056. Если цифровая
грамотность оказывает положительное влияние на ВВП, то ее значимость
составляет 10\%. Это связано с важностью цифрового образования, которое
в долгосрочной перспективе может способствовать повышению
производительности и инновационного потенциала.

Индекс человеческого капитала (\emph{Х4}) - \emph{F-Value} = 1,85 и
\emph{P-Value} =0,207. Согласно данной модели, индекс человеческого
капитала не оказывает статистически значимого влияния на ВВП. Это может
указывать на то, что улучшение человеческого капитала (образование,
здравоохранение, навыки) не оказывает прямого влияния на рост ВВП,
особенно в краткосрочной перспективе.

Производительность труда (\emph{Х5}). \emph{F-Value} = 7,13 и
\emph{P-Value} =0,026. Производительность труда в секторе "информация и
связь" оказывает значительное негативное влияние на ВВП Казахстана. Это
связано с тем, что повышение производительности в некоторых секторах
требует значительных инвестиций в технологии, которые могут не оказывать
ощутимого прямого воздействия.

Доля населения, пользующегося интернетом (\%) (\emph{Х6}) -
\emph{F-Value} = 7,31 и \emph{P-Value} =0,024. Процент пользователей
Интернета оказывает значительное негативное влияние на ВВП, что может
быть связано с длительным процессом адаптации населения к новым цифровым
технологиям. Этот фактор может отражать разницу между доступностью и
качеством Интернета, что также влияет на экономические последствия.

\emph{Качество модели}.

Ошибка (\emph{Error}): средняя ошибка модели (среднеквадратичная ошибка)
составляет 516,9, что указывает на определенную ошибку прогнозирования,
но незначительную с точки зрения общей дисперсии.

Общая вариация (\emph{Total SS}=33751,8) и объясненная вариация
(\emph{Regressian SS} =29 100,0) подтверждает высокую подгонку модели с
\emph{R²} = 86,22\%. Другими словами, изменение номинального ВВП
Казахстана на 86\% можно объяснить выбранными факторами. Однако
показатель \emph{R²} (pred) = 0,00\% свидетельствует о том, что модель
плохо предсказывает значения ВВП за пределами исследуемого периода.

На основе полученной модели выполнен прогноз номинального ВВП на
2026-2028гг. на основе методики наименьших квадратов (OLS). На основе
оцененной модели рассчитывается номинальный прогноз ВВП для трехлетнего
периода (таблица 6).

Прогноз показывает стабильную положительную динамику. В соответствии с
текущей тенденцией цифровой трансформации экономики, ожидается, что
номинальный ВВП вырастет на 18--19\% в период с 2026 по 2028 годы.
Результирующий путь соответствует сценарию безынерционного
экономического развития при сохранении текущих темпов цифровой
трансформации Казахстана. Рост ВВП объясняется совокупным влиянием
цифровых инвестиций, корпоративных и производственных факторов.

\tcap{Таблица 6 -- Краткосрочный прогноз}
\begin{longtblr}[
  label = none,
  entry = none,
]{
  cells = {c},
  cells = {font = \small},
  row{1} = {font=\bfseries\small},
  hlines,
  vlines,
}
Факторы & 2026 & 2027 & 2028\\
У & 320,26 & 348,86 & 379,97\\
Х1 & 332,79 & 374,42 & 418,79\\
Х2 & 1 422,15 & 1 600,70 & 1 791,05\\
Х3 & 96,05 & 98,11 & 100,17\\
Х4 & 0,67 & 0,63 & 0,59\\
Х5 & 150 620,13 & 162 034,70 & 174 042,00\\
Х6 & 42,11 & 29,62 & 15,62
\end{longtblr}

\begin{multicols}{2}
Прогноз показывает устойчивый рост государственных средств на цифровые
технологии. Среднегодовой темп роста составляет более 10\%, что
свидетельствует о стратегическом приоритете цифровизации в
государственной политике РК. Увеличение государственных инвестиций
создаст инфраструктурную базу для расширения цифровых услуг и
модернизации экономического развития.

Частный сектор демонстрирует сопоставимую динамику роста инвестиционных
вложений. Увеличение инвестиций отражает рост технологической адаптации
предприятий и компаний, а также хорошее укрепление конкурентной среды.
Синергетический эффект государственных и частных капиталовложений
создает определенный стабильный мультипликативный эффект для
национального хозяйства.

Таким образом, согласно данному прогнозу, уровень обеспеченности граждан
Казахстана цифровыми компетенциями приближается к максимальному.
Значения, близкие к 100\%, свидетельствуют о высокой степени цифровой
интеграции общества. Дальнейшая динамика данного показателя, вероятно,
будет носить преимущественно качественный и стабилизационный характер.

Умеренная вариативность индекса человеческого капитала отражает
структурные изменения в образовании и социально-экономической среде.
Показатель остается в диапазоне, соответствующем уровню устойчивого
человеческого развития. Следует иметь в виду, что этот индекс является
интегральным и чувствительным к различным макроэкономическим факторам.

Наблюдается устойчивый рост производительности труда, что подтверждает
повышение продуктивности цифрового сектора. Увеличение этого показателя
является прямым каналом влияния цифровизации на формирование повышенной
ценности. Производительность и рост рабочей силы увеличивают вклад
информационно-коммуникационного сектора в структуру ВВП.

Динамика показателя доли граждан, использующего интернет, относится к
производственным характеристикам моделей временных рядов. В долгосрочной
перспективе доступ к Интернету остается на высоком уровне. В рамках
модели этот показатель интерпретируется как структурный фактор цифровой
трансформации, связанный с экономическим потенциалом за счет расширения
цифровых рынков и услуг.

Прогнозная динамика показывает улучшения, связанные с ключевыми цифрами
и макроэкономическими показателями. Увеличение инвестиций в цифровые
технологии является ресурсной основой для повышения производительности
труда. Рост цифровых навыков населения создает благоприятную
институциональную среду для расширения цифровой экономики. В целом, эти
факторы подтверждают положительную тенденцию повышения номинального ВВП
Казахстана в среднесрочной перспективе.

Следовательно, результаты моделирования и прогнозирования подтвердили
приведенную гипотезу, где независимые переменные - инвестиции
государства и предприятий в электронные и цифровые технологии -
оказывают достаточно сильное влияние на степень цифровизации управления
национальной экономики Казахстана. Значит, эти финальные результаты
могут оказать помощь для принятия своевременных и обоснованных решений
государства в ближайшей перспективе.

{\bfseries Выводы.} Реализация эконометрического моделирования за период
2010--2025 годов показала, что в общем, на экономический рост оказывает
влияние рост инвестиций в цифровые технологии и повышение
производительности труда в Казахстане. Вместе с тем цифровая грамотность
непосредственно усиливает это воздействие, а показатели распространения
интернета и уровень человеческого капитала оказывают менее выраженное
краткосрочное влияние.

\href{https://www.tandfonline.com/author/la+Bella\%2C+Marco+V+L}{Bella},
\href{https://www.tandfonline.com/author/Santoro\%2C+Patrizia}{Santoro}
(2025) справедливо утверждают, что необходимо усиленно развивать
внедрение электронных государственных услуг {[}27{]}. Проведённый анализ
показал, что необходимо осуществить трансформацию органов власти с
позиции усиления использования цифровых технологий и повышения
квалификации государственных служащих.

В июне 2025 г. для приоритетного развития свыше 400 госуслуг
инвестиционным проектам был введен новый Механизм «Зеленого коридора» на
базе информационного портала «Е-лицензирование». Данный механизм
применяется при сотрудничестве инвесторов с государственными органами в
достаточно ускоренные сроки на уровне «одного окна».

Такая процедура предоставляет возможность оперативно сопровождать заявки
на инвестиционные проекты в рамках использования цифровых
государственных сервисов на основе НЦИП (Национальной цифровой
инвестиционной платформы). Такая практика стимулирует организацию
гарантированного и последовательного порядка рассмотрения инвестиционных
проектов на базе оказания государственных услуг.

С помощью пространственно-временных индикаторов Bobro et al. (2025)
установили главные векторы инновационного развития процесса управления
Европы на основе ИКТ {[}28{]}. Здесь особый акцент авторы сделали на
предоставлении качественных государственных услуг. В результате было
определено, что инструменты искусственного интеллекта способствуют
продвижению роста эффективности государственного управления.

Проведенный анализ показал, что уровень цифровой граммотности
казахстанского населения с 2010 года до 2025 года возрос с 62\% до 94\%.
В настоящее время электронная платформа «eOtinish» - стала важным
подходом на уровне цифрового взаимодействия граждан с государством. К
примеру, если в 2022 году количество обращений населения составляло 1.8
млн. единиц, то по результатам 2025 года -- более 5 млн. ед. (рост почти
за 4 года составил в 2,8 раза). В Казахстане платформа «eOtinish» была
введена летом 2022г., где внедрена автоматизированная процедура подачи,
отбора и мониторинга обращений, что позволяет пользователям оперативно
отслеживать статус обращений и заявлений в государственные органы в
онлайн режиме, что соответствует использованию на практике принципов
«Слышащего государства». Это свидетельствует о том, что платформа
«eOtinish» оказывает положительное влияние и способствует развитию
прозрачного, доверительного и результативного взаимодействия государства
с гражданами Казахстана в цифровой среде.

Вместе с тем, в ближайшей перспективе созданы предпосылки для внедрения
методов ИИ и ускорения процессов взаимодействия госорганов с населением
в мобильных приложениях и цифровых порталах.

В Казахстане с 2025 года стартовала программа по обучению 50.000
госслужащих по работе с новыми информационными программами с применением
ИИ. Следует отметить, что для продвижения «IT-стартапов», роста
профессионализма и борьбы с коррупцией в системе государственной службы
активизированы мероприятия по повышению квалификации государственных
работников в области приобретения навыков по ИИ. По данным Министерства
ИИ и цифрового развития на 1-ый квартал 2026 года свыше 3000 госслужащих
уже успешно прошли обучение по обретению профессиональных навыков в
работе с инструментами ИИ. 

Итальянскими авторами (\href{javascript:;}{Todisco} et al., 2025) был
разработан и реализован опрос в онлайн формате среди 42 государственных
работников, чтобы оценить как метод «anticipatory governance» (AG)
воздействует на их уровень инновационного развития {[}29{]}. Здесь было
обоснованно, что методика «anticipatory governance» (AG) включает
владение госслужащими цифровыми навыками, которые стимулируют принятие
объективных и справедливых упреждающих государственных решений.

Становление архитектуры ИИ уже начало позитивно воздействовать на
деятельность определенных отраслей. В частности, в секторе энергетики
применение методики ИИ помогает воздействовать на увеличение
энергоэффективности и перераспределения энергомощностей в пиковые часы.

Определенная поддержка использования подходов ИИ происходит в сфере
транспорта и логистики, что привело к возрастанию на 10 процентов
количества грузоперевозок, снижения логистических расходов до 20\% в
2025 году.

Например, казахстанский проект «Социальный кошелек» включил около 7000
образовательных школ и 49 млн обедов на бесплатной основе. Это дало
возможность снизить расходы и сэкономить средства бюджета до 30\%.
В государственных аптеках внедрение цифровизации и ИИ-инструментов
включило 13 млн. электронных лекарственных рецептов.

Однако, в республике еще имеются значительные угрозы и барьеры развития
цифровой трансформации: низкий уровень доступа интернета и развитости
информационной инфраструктуры в отдельных районах и сельской местности.

Несмотря на запуск 2-х суперкомпьютеров в Казахстане, наблюдается
серьезный разрыв с международными ведущими IT-компаниям. Так,
у международных IT-гигантов производительный потенциал
составляет несколько миллионов видеокарт, что позволяет им располагать
высокими мощностями. В нашей стране таких значительных технологических
ресурсов недостаточно еще!

Очевидно, что информационная база государственных услуг требует еще
доработки, существуют еще перебои в деятельности электронных платформ
в ЦОНах.

Внедрение ИИ имеет сегодня еще такие ограничения, как: высокая цена
приобретения инновационного оборудования, довольно большая стоимость
вычислений, низкая степень развития «5G сетей», определенная и слабая
доступность «Graphics Processing Unit» (GPU -- то есть - это главный
вычислительный чип видеокарты) из-за дорогих издержек для
предпринимательских структур.

Статья Bastykov et al. (2024) посвящена реформам в системе
государственного регулирования экономики. Здесь были охарактеризованы
основные направления по осуществлению этапов и рекомендаций по
стимулированию развития электронного правительства {[}30{]}.

Поэтому с целью совершенствования и продвижения системы ИИ и внедрения
ИКТ в государственном регулировании экономики РК, нами предлагаются
\emph{такие предложения}.

Госорганам следует постоянно мониторить персональную ответственность
пользователей ИИ-систем, чтобы поддерживать безопасность.

Использовать риск-ориентированный подход при сборе и анализе статданных
на основе инновационных технологий.

Расширять приток инвестиций в технологическое развитие ключевых секторов
экономики, как со стороны Правительства, так и бизнеса.

Увеличивать производственную мощность информационной и сетевой
инфраструктуры ИИ от киберугроз, разного рода вирусов и пр.

Создавать мощные дата-центры для размещения серверного оборудования,
деятельности интернет-сервисов и облачных платформ, хранения крупных баз
данных с позиции достижения экономической безопасности, обеспечения
высокоскоростными интернет-каналами, оптимального управления в
госорганах, эффективного использования производственных процессов для
бизнес-структур и скорости обработки данных населения.

В целом, с 2025 г. ИИ стал основным инструментом при принятии
государственных решений в экономике на базе введения «Национальной
платформы ИИ». Посредством данной платформы все госорганы имеют доступ
ко всем имеющимся вычислительным средствам, электронным сервисам, единым
статданным, подходам, IT-проектам и моделям.

Таким образом, принятие соответствующих нормативных документов по ИКТ и
ИИ, создание сети новых цифровых платформ, расширение инвестиций
государства и предприятий в ИИ, активизация проникновения интернета,
повышение квалификации госслужащих, достаточно развитая экосистема
государственных электронных услуг - представляют благоприятную основу
для осуществления перспектив цифровой трансформации, внедрения
инструментов ИИ, работу дата-центров в управлении и достижении
устойчивого экономического развития Казахстана.
\end{multicols}

\begin{center}
{\bfseries Литература}
\end{center}

\begin{refs}
1. Guterres, А. From India, Guterres calls for \$3 billion fund to ensure
AI benefits all. - United Nations, 2026. URL:
\url{https://news.un.org/en/story/2026/02/1166996.-} Дата обращения:
22.02.2026.

2. Замедление роста, цифровизация и тарифы: что ждет мировую торговлю в
2026 году? \href{https://news.un.org/ru/}{Новости ООН}. Глобальный
взгляд Человеческие судьбы, 2026. URL:
https://news.un.org/ru/story/2026/01/1467186.- Дата обращения:
28.01.2026.

3. Sigurjonsson, О., Jónsson, Е., Gudmundsdottir, S. Sustainability of
Digital Initiatives in Public Services in Digital Transformation of
Local Government: Insights and Implications // Sustainability. - 2024.
- Vol.16(24):10827. DOI 10.3390/su162410827.

4. \href{https://datareportal.com/reports/digital-2026-global-overview-report}{Digital
2026 Global Overview Report} / 2025. URL:
\url{https://wearesocial.com/us/blog/2025/10/digital-2026-global-overview-report/}
.-Дата обращения: 23.01.2026.

5. Smagulova Sh., Bastykov D., Omarova А., Omarova В. The evolution and
development of public administration digitalization: аn international
bibliometric analysis // Cogent Social Sciences. - 2025. - Vol.11(1):
2548018. DOI~10.1080/23311886.2025.2548018.

6. Цифровой Кодекс РК / Кодекс РК от 9 января 2026 г. - № 255-VIII ЗРК.
-- Астана: АкОрда, 2026. URL:
\url{https://adilet.zan.kz/rus/docs/K2600000255.-} Дата обращения:
19.01.2026.

7. Об искусственном интеллекте / Закон РК от 17 ноября 2025 г. - №
230-VIII ЗРК. - Астана: АкОрда, 2025. URL:
https://adilet.zan.kz/rus/docs/Z2500000230.- Дата обращения:
20.02.2026.

8. Концепция развития искусственного интеллекта на 2024--2029 гг. /
Постановление Правительства РК от 24 июля 2024 г. - № 592. -- Астана:
АкОрда, 2024. URL:
\url{https://continent-online.com/Document/?doc_id=34317430}.- Дата
обращения 25.02.2026.

9. Tutar H. Algorithmic Interaction in Public Administration: Data-Based
Communication and Political Participation in Public Institutions //
Public Governance, Administration and Finances Law Review. - 2025. -
Vol.10 (2). - Р.157-184. DOI
\href{https://doi.org/10.53116/pgaflr.8460}{10.53116/pgaflr.8460}.

10. \href{https://journals.sagepub.com/doi/10.1177/00208523251333541\#con}{Borriello}
G. The grand challenge of public administration digitalization: The
digital identity policy in Italy \emph{//}
\href{https://journals.sagepub.com/home/RAS}{International Review of
Administrative Sciences}.-2025. - Vol.2. Р.281-319. DOI
10.1177/00208523251333541.

11. Koerich A.B., Mussi C.C., Salgueirinho J.B., Guerra O.A., Casagrande
J.L. Critical Factors in the Implementation of Innovation: A Case
Study of a Brazilian Public Organization // The Electronic Journal of
Information Systems in Developing Countries. - 2025. - Vol.4. -
Р.1-15. DOI 10.1002/isd2.70018.

12. Yang D., Jun X., Xiao Y.
\href{https://ideas.repec.org/a/eee/teinso/v83y2025ics0160791x25001915.html}{How
public digital governance system affects firms'{}
digital technology innovation performance: Base on open innovation
perspective} //
\href{https://ideas.repec.org/s/eee/teinso.html}{Technology in
Society}.- 2025. - Vol.83: 103001. DOI 10.1016/j.techsoc.2025.103001.

13. Казахстан лидирует по готовности к внедрению искусственного интеллекта
среди стран Центральной Азии / Министерство искусственного интеллекта
и цифрового развития. - Астана, 2025. URL:
\url{https://www.gov.kz/memleket/entities/maidd/press/}
news/details/1131135?lang=ru. -Дата обращения 10.01.2026.

14. Chen J., Lou~Z., Chen G.,~ Zhang~C.,~Li Y. Research on the Application
of AI-Based Intelligent Decision Support Systems in E-Government
Services //
\href{https://www.sciencedirect.com/org/journal/journal-of-organizational-and-end-user-computing}{Journal
of Organizational and End User Computing}.-2025. - Vol.37(1). -
Р.1-25. DOI 10.4018/JOEUC.392618.

15. Казахстан вошел в топ-10 в мире по уровню развития цифровых услуг /
Международное информационное агентство «Казинформ», 2025. URL:
\url{https://www.inform.kz/ru/kazahstan-voshel-v-top-10-v-mire-po-urovnyu-razvitiya-tsifrovih-uslug-f606f5}
. - Дата обращения 08.01.2026.

16. Statistics of services rendered / Portal «egov». -- Astana, 2026. URL:
https://egov.kz/cms/en/information/about/stat.- Дата обращения:
20.01.2026.

17. Национальный доклад о состоянии государственной службы. -- Aстана,
2025. URL:
\url{file:///C:/Users/777/Downloads/_gluster_2025_4_10_105485a81d162de727bef20a077e631f_original.628447.pdf}
.- Дата обращения: 30.12.2025.

18. Wages in the Republic of Kazakhstan / Agency for Strategic Planning
and Reforms of the RK. Bureau of National Statistics. -- Astana, 2026.
URL:
\url{https://stat.gov.kz/en/industries/labor-and-income/stat-wags/publications/184222/?utm_source=chatgpt.com}
.- Дата обращения: 20.02.2026.

19. Dei, Н. The use of AI in the organization of local government work //
LatIA. -2025. - Vol.3. - Р.123-133. DOI
\href{https://doi.org/10.62486/latia2025123}{10.62486/latia2025123}.

20. \href{https://www.researchgate.net/scientific-contributions/Darusalam-Darusalam-2146395918?_tp=eyJjb250ZXh0Ijp7ImZpcnN0UGFnZSI6ImNpdGF0aW9uRG93bmxvYWQiLCJwYWdlIjoicHVibGljYXRpb24ifX0}{Darusalam}
D.,
\href{https://www.researchgate.net/profile/Marijn-Janssen?_tp=eyJjb250ZXh0Ijp7ImZpcnN0UGFnZSI6ImNpdGF0aW9uRG93bmxvYWQiLCJwYWdlIjoicHVibGljYXRpb24ifX0}{Janssen}
M.,
\href{https://www.researchgate.net/scientific-contributions/Sri-Delasmi-Jayanti-2236563283}{Jayanti}
S.,
\href{https://www.researchgate.net/scientific-contributions/Renato-Sitompul-2229051296?_tp=eyJjb250ZXh0Ijp7ImZpcnN0UGFnZSI6ImNpdGF0aW9uRG93bmxvYWQiLCJwYWdlIjoicHVibGljYXRpb24ifX0}{Sitompul}
R. Public administration digitalization effects on corruption: Lesson
learned from Indonesia // Digital Government: Research and Practice.-
2024. - Vol.5(4). - Р.1-13. DOI
\href{https://doi.org/10.1145/3691351}{10.1145/3691351}.

21. Urhibo B., Aruoren Е., Ivwighren Н., Imene О. Digital Communication
Technologies in Public Administration: A Study of Ministry of Basic
and Secondary Education, Delta State//
\href{https://www.researchgate.net/journal/Journal-of-Ecohumanism-2752-6801?_tp=eyJjb250ZXh0Ijp7ImZpcnN0UGFnZSI6ImNpdGF0aW9uRG93bmxvYWQiLCJwYWdlIjoicHVibGljYXRpb24ifX0}{Journal
of Ecohumanism}.- 2024. - Vol.3(7). - Р.319--328. DOI
\href{https://doi.org/10.62754/joe.v3i7.4203}{10.62754/joe.
v3i7.4203}.

22. GDP (current US\$) / World Bank. -- Washington DC, 2025. URL:
https://data.worldbank.org/indicator/NY.GDP.MKTP.CD?locations=KZ\&utm\_source=chatgpt.com.-
Дата обращения: 20.01.2026.

23. State program "Digital Kazakhstan" / Order of the Government of the
Republic of Kazakhstan. - Astana, 2017. URL:
\url{https://cis-legislation.com/document.fwx?rgn=102961}
\&utm\_source=chatgpt.com. -Аccessed on 30.12.2025.

24. Об утверждении Концепции цифровой трансформации, развития отрасли
информационно-коммуникационных технологий и кибербезопасности на
2023--2029 годы/Постановление Правительства Республики
Казахстан от 28 марта 2023 года № 269. -- Aстана: Акорда, 2023. URL:
\url{https://stat.gov.kz/en/industries/labor-and-income/stat-wags/publications/184222/?utm_source=chatgpt.com.-}
Дата обращения 01.15.2026.

25. Artificial Intelligence Is the New Literacy: Kazakhstan Prepares Its
People for a Digital Future / DKN. World news. - Astana, 2026. URL:
https://dknews.kz/en/articles-in-english/380871-artificial-intelligence-is-the-new-literacy?utm\_source=chatgpt.com.-Дата
обращения: 09.01.2026.

26. Kazakhstan internet users / Trendonify. -- Astana, 2025. URL:
https://trendonify.com/kazakhstan/internet-users?utm\_source=chatgpt.com.-
Дата обращения: 05.01.2026.

27. \href{https://www.tandfonline.com/author/la+Bella\%2C+Marco+V+L}{Bella}
М.,
\href{https://www.tandfonline.com/author/Santoro\%2C+Patrizia}{Santoro}
Р. The Modernization of Public Administration: The Impact of Public
Employment on Digital Innovation in Italy //
\href{https://www.tandfonline.com/journals/lpad20}{International
Journal of Public Administration~}Volume\emph{.} - 2025. - Vol.48(4).
- Р.250-262. DOI 10.1080/01900692.2024.2346784.

28. Bobro N., Bielikov V., Matveyeva М., Salamakha А., Kharchun, V.
Innovations in public administration and management for implementing
effective strategies and tools // Salud, Ciencia y Tecnología -- Serie
de Conferencias. -2025. - Vol.4. -Р.1-9. DOI
10.56294/sctconf20251483.

29. \href{javascript:;}{Todisco} L., \href{javascript:;}{Mangia} G.,
\href{javascript:;}{Canonico} P., \href{javascript:;}{Langella} F.
Foreseeing the future: anticipatory governance as a response to the
technological and managerial challenges in the public sector //
International Journal of Public Sector Management.-2025. - Vol.3. -
Р.1-17. DOI 10.1108/IJPSM-03-2025-0109.

30. Bastykov, D., Smagulova, Sh., Karatayev, M., Тogizbaeva, M.
Bibliometric analysis of research on public administration reform //
Journal of Economic Research \&Amp; Business Administration. - 2024. -
Vol.149(3). - Р.136-149.
DOI~\href{https://doi.org/10.26577/be.2024-149-i3-010}{10.26577/be.2024-149-i3-010}.
\end{refs}

\begin{center}
{\bfseries References}
\end{center}

\begin{refs}
1. Guterres, А. From India, Guterres calls for \$3 billion fund to ensure
AI benefits all. - United Nations, 2026. URL:
\url{https://news.un.org/en/story/2026/02/1166996.-} Аccessed on
22.02.2026.

2. Zamedlenie rosta, cifrovizacija i tarify: chto zhdet mirovuju
torgovlju v 2026 godu? Novosti OON. Global' nyj vzgljad
Chelovecheskie sud' by, 2026. URL:
https://news.un.org/ru/story/2026/01/1467186.- Data obrashhenija:
28.01.2026. {[}in Russian{]}

3. Sigurjonsson, О., Jónsson, Е., Gudmundsdottir, S. Sustainability of
Digital Initiatives in Public Services in Digital Transformation of
Local Government: Insights and Implications // Sustainability. - 2024. -
Vol.16(24):10827. DOI 10.3390/su162410827.

4. \href{https://datareportal.com/reports/digital-2026-global-overview-report}{Digital
2026 Global Overview Report} / 2025. URL:
\url{https://wearesocial.com/us/blog/2025/10/digital-2026-global-overview-report/}-
Аccessed on 23.01.2026.

5. Smagulova Sh., Bastykov D., Omarova А., Omarova В. The evolution and
development of public administration digitalization: аn international
bibliometric analysis // Cogent Social Sciences. - 2025. - Vol.11(1):
2548018. DOI~10.1080/23311886.2025.2548018.

6. Cifrovoj Kodeks RK / Kodeks RK ot 9 janvarja 2026 g. - № 255-VIII ZRK.
-- Astana: AkOrda, 2026. URL:
https://adilet.zan.kz/rus/docs/K2600000255.- Data obrashhenija:
19.01.2026. {[}in Russian{]}

7. Ob iskusstvennom intellekte / Zakon RK ot 17 nojabrja 2025 g. - №
230-VIII ZRK. - Astana: AkOrda, 2025. URL:
https://adilet.zan.kz/rus/docs/Z2500000230.- Data obrashhenija:
20.02.2026. {[}in Russian{]}

8. Koncepcija razvitija iskusstvennogo intellekta na 2024--2029 gg. /
Postanovlenie Pravitel' stva RK ot 24 ijulja 2024 g. - №
592. -- Astana: AkOrda, 2024. URL:
https://continent-online.com/Document/?doc\_id=34317430.- Data
obrashhenija 25.02.2026. {[}in Russian{]}

9. Tutar H. Algorithmic Interaction in Public Administration: Data-Based
Communication and Political Participation in Public Institutions //
Public Governance, Administration and Finances Law Review. - 2025. -
Vol.10 (2). - Р.157-184. DOI
\href{https://doi.org/10.53116/pgaflr.8460}{10.53116/pgaflr.8460}.

10. \href{https://journals.sagepub.com/doi/10.1177/00208523251333541\#con}{Borriello}
G. The grand challenge of public administration digitalization: The
digital identity policy in Italy \emph{//}
\href{https://journals.sagepub.com/home/RAS}{International Review of
Administrative Sciences}.-2025. - Vol.2. Р.281-319. DOI
10.1177/00208523251333541.

11. Koerich A.B., Mussi C.C., Salgueirinho J.B., Guerra O.A., Casagrande
J.L. Critical Factors in the Implementation of Innovation: A Case Study
of a Brazilian Public Organization // The Electronic Journal of
Information Systems in Developing Countries. - 2025. - Vol.4. - Р.1-15.
DOI 10.1002/isd2.70018.

12. Yang D., Jun X., Xiao Y.
\href{https://ideas.repec.org/a/eee/teinso/v83y2025ics0160791x25001915.html}{How
public digital governance system affects firms'{} digital
technology innovation performance: Base on open innovation perspective}
// \href{https://ideas.repec.org/s/eee/teinso.html}{Technology in
Society}.- 2025. - Vol.83: 103001. DOI 10.1016/j.techsoc.2025.103001.

13. Kazahstan lidiruet po gotovnosti k vnedreniju iskusstvennogo
intellekta sredi stran Central' noj Azii / Ministerstvo
iskusstvennogo intellekta i cifrovogo razvitija. - Astana, 2025. URL:
https://www.gov.kz/memleket/entities/maidd/press/
news/details/1131135?lang=ru. -Data obrashhenija 10.01.2026.

14. Chen J., Lou~Z., Chen G.,~ Zhang~C.,~Li Y. Research on the
Application of AI-Based Intelligent Decision Support Systems in
E-Government Services //
\href{https://www.sciencedirect.com/org/journal/journal-of-organizational-and-end-user-computing}{Journal
of Organizational and End User Computing}.-2025. - Vol.37(1). - Р.1-25.
DOI 10.4018/JOEUC.392618.

15. Kazahstan voshel v top 10 v mire po urovnju razvitija cifrovyh uslug
/ Mezhdunarodnoe informacionnoe agentstvo «Kazinform», 2025. URL:
https://www.inform.kz/ru/kazahstan-voshel-v-top-10-v-mire-po-urovnyu-razvitiya-tsifrovih-uslug-f606f5.
- Data obrashhenija 08.01.2026. {[}in Russian{]}

16. Statistics of services rendered / Portal «egov». -- Astana, 2026.
URL: https://egov.kz/cms/en/information/about/stat.- Аccessed on
20.01.2026.

17. Nacional' nyj doklad o sostojanii gosudarstvennoj
sluzhby. -- Astana, 2025. URL:
\_gluster\_2025\_4\_10\_105485a81d162de727bef20a077e631f\_original.628447.pdf .-
Data obrashhenija: 30.12.2025.

18. Wages in the Republic of Kazakhstan / Agency for Strategic Planning
and Reforms of the RK. Bureau of National Statistics. -- Astana, 2026.
URL:
\url{https://stat.gov.kz/en/industries/labor-and-income/stat-wags/publications/184222/?utm_source=chatgpt.com}.-
Аccessed on 20.02.2026.

19. Dei, Н. The use of AI in the organization of local government work //
LatIA. -2025. - Vol.3. - Р.123-133. DOI
\href{https://doi.org/10.62486/latia2025123}{10.62486/latia2025123}.

20. \href{https://www.researchgate.net/scientific-contributions/Darusalam-Darusalam-2146395918?_tp=eyJjb250ZXh0Ijp7ImZpcnN0UGFnZSI6ImNpdGF0aW9uRG93bmxvYWQiLCJwYWdlIjoicHVibGljYXRpb24ifX0}{Darusalam}
D.,
\href{https://www.researchgate.net/profile/Marijn-Janssen?_tp=eyJjb250ZXh0Ijp7ImZpcnN0UGFnZSI6ImNpdGF0aW9uRG93bmxvYWQiLCJwYWdlIjoicHVibGljYXRpb24ifX0}{Janssen}
M.,
\href{https://www.researchgate.net/scientific-contributions/Sri-Delasmi-Jayanti-2236563283}{Jayanti}
S.,
\href{https://www.researchgate.net/scientific-contributions/Renato-Sitompul-2229051296?_tp=eyJjb250ZXh0Ijp7ImZpcnN0UGFnZSI6ImNpdGF0aW9uRG93bmxvYWQiLCJwYWdlIjoicHVibGljYXRpb24ifX0}{Sitompul}
R. Public administration digitalization effects on corruption: Lesson
learned from Indonesia // Digital Government: Research and Practice.-
2024. - Vol.5(4). - Р.1-13. DOI
\href{https://doi.org/10.1145/3691351}{10.1145/3691351}.

21. Urhibo B., Aruoren Е., Ivwighren Н., Imene О. Digital Communication
Technologies in Public Administration: A Study of Ministry of Basic and
Secondary Education, Delta State//
\href{https://www.researchgate.net/journal/Journal-of-Ecohumanism-2752-6801?_tp=eyJjb250ZXh0Ijp7ImZpcnN0UGFnZSI6ImNpdGF0aW9uRG93bmxvYWQiLCJwYWdlIjoicHVibGljYXRpb24ifX0}{Journal
of Ecohumanism}.- 2024. - Vol.3(7). - Р.319--328. DOI
\href{https://doi.org/10.62754/joe.v3i7.4203}{10.62754/joe. v3i7.4203}.

22. GDP (current US\$) / World Bank. -- Washington DC, 2025. URL:
https://data.worldbank.org/indicator/NY.GDP.MKTP.CD?locations=KZ\&utm\_source=chatgpt.com.-
Аccessed on 20.01.2026.

23. State program "Digital Kazakhstan" / Order of the Government of the
Republic of Kazakhstan. - Astana, 2017. URL:
\url{https://cis-legislation.com/document.fwx?rgn=102961}
\&utm\_source=chatgpt.com. - Аccessed on 30.12.2025

24. Ob utverzhdenii Koncepcii cifrovoj transformacii, razvitija otrasli
informacionno-kommunikacionnyh tehnologij i kiberbezopasnosti na
2023--2029 gody/Postanovlenie Pravitel' stva Respubliki
Kazahstan ot 28 marta 2023 goda № 269. -- Astana: Akorda, 2023. URL:
https://stat.gov.kz/en/industries/labor-and-income/stat-wags/publications/184222/?utm\_source=chatgpt.com.-Data
obrashhenija 01.15.2026.

25. Artificial Intelligence Is the New Literacy: Kazakhstan Prepares Its
People for a Digital Future / DKN. World news. - Astana, 2026. URL:
https://dknews.kz/en/articles-in-english/380871-artificial-intelligence-is-the-new-literacy?utm\_source=chatgpt.com.-
Аccessed on 09.01.2026.

26. Kazakhstan internet users / Trendonify. -- Astana, 2025. URL:
\url{https://trendonify.com/kazakhstan/internet-users?utm_source=chatgpt.com.-}
Аccessed on 05.01.2026.

27. \href{https://www.tandfonline.com/author/la+Bella\%2C+Marco+V+L}{Bella}
М.,
\href{https://www.tandfonline.com/author/Santoro\%2C+Patrizia}{Santoro}
Р. The Modernization of Public Administration: The Impact of Public
Employment on Digital Innovation in Italy //
\href{https://www.tandfonline.com/journals/lpad20}{International Journal
of Public Administration~}Volume\emph{.} - 2025. - Vol.48(4). - Р.
250-262. DOI 10.1080/01900692.2024.2346784.

28. Bobro N., Bielikov V., Matveyeva М., Salamakha А., Kharchun, V.
Innovations in public administration and management for implementing
effective strategies and tools // Salud, Ciencia y Tecnología -- Serie
de Conferencias. -2025. - Vol.4. -Р.1-9. DOI 10.56294/sctconf20251483.

29. \href{javascript:;}{Todisco} L., \href{javascript:;}{Mangia} G.,
\href{javascript:;}{Canonico} P., \href{javascript:;}{Langella} F.
Foreseeing the future: anticipatory governance as a response to the
technological and managerial challenges in the public sector //
International Journal of Public Sector Management.-2025. - Vol.3. -
Р.1-17. DOI 10.1108/IJPSM-03-2025-0109.

30. Bastykov, D., Smagulova, Sh., Karatayev, M., Тogizbaeva, M.
Bibliometric analysis of research on public administration reform //
Journal of Economic Research \&Amp; Business Administration. - 2024. -
Vol.149(3). - Р.136-149.
DOI~\href{https://doi.org/10.26577/be.2024-149-i3-010}{10.26577/be.2024-149-i3-010}
\end{refs}

\begin{info}
\emph{{\bfseries Сведения об авторах}}

Смагулова Ш.А. - доктор экономических наук, профессор, Университет
Международного Бизнеса им. Кенжегали Сагадиева, Алматы, Казахстан,
е-mail: shsmagulova@mail.ru;

Кажимуратова С.З. - докторант PhD, НАО Университет Нархоз, Алматы,
Казахстан, е-mail: saltanat.kazhimuratova@narxoz.kz;

Койшинова Г.К. - кандидат экономических наук, ассоциированный
профессор, ВКУ им. С.Аманжолова, Усть-Каменогорск, Казахстан, е-mail:
gaini\_1964@ mail.ru;

Жетписбаева А.А. - старший преподаватель, Казахский Национальный
педагогический университета им. Абая, Алматы, Казахстан, е-mail:
zhetpisbayeva.asem@inbox.ru;

Жакупова С.Т. - кандидат экономических наук, доцент, Университет
Международного Бизнеса им. Кенжегали Сагадиева, Алматы, Казахстан,
е-mail: Saltanat.zhakupova.69@gmail.ru;

Ахметова М.Д. - магистр экономических наук, Университет Нархоз,
Алматы, Казахстан, е-mail: kz1985@mail.ru;

\emph{{\bfseries Information about the authors}}

Smagulova Sh. - Doctor of Economics, Professor, Kenzhegali Sagadiyev
University of International Business, Almaty, Kazakhstan, е-mail:
shsmagulova@mail.ru;

Kazhimuratova S.- PhD student ,Narxoz University, Almaty, Kazakhstan,
е-mail: saltanat.kazhimuratova@narxoz.kz;

Koishinova G.K.- PhD in Economics, Associate Professor, S. Amanzholov
East Kazakhstan University,

Ust-Kamenogorsk, Kazakhstan, е-mail: gaini\_1964@mail.ru;

Zhetpisbayeva A. - Senior Lecturer, Abay Kazakh National Pedagogical
University, Almaty, Kazakhstan, е-mail: zhetpisbayeva.asem@ibbox.ru;

Zhakupova S. - Candidate of Economic Sciences, Associate Professor,
Kenzhegali Sagadiyev University of International Business, Almaty,
Kazakhstan, е-mail: Saltanat.zhakupova.69@gmail.ru.

Akhmetova M. master of economic science, Narxoz University, Almaty,
Kazakhstan, е-mail: kz1985@mail.ru.
\end{info}
